% Introduction. Architecture follows CH7 (question, taxonomy, results in
% advance, organisation) and CH2 (threads), not the four-preface CH8 template.
% Streamlined after the removal of the two-phase section and the demotion of
% the birth--plague and public-good models to appendices: the taxonomy figure
% is gone and its content carried in prose, the headline results are three
% rather than four, and the notation table covers the body only.

\section{Introduction}
\label{var:sec:introduction}

Every model carried through to a result in the preceding chapters holds its
rates fixed. A replicative unit is born at rate $\lambda$, dies at rate $\mu$
and triggers catastrophe at rate $\delta$, and those three numbers are the
same at the end of the process as at the beginning. That assumption is what
makes the generating functions close, and it buys a great deal: the roots $a$
and $b$, the burst-size law, the flooding criterion and the removal budget.
It is also the first thing an ecologist will question, because in almost
every setting one might want to describe, a rate is not a constant ---
resources run out, competitors adapt, and success carries its own cost.

This chapter asks what happens when a rate is allowed to move. The models
collected here were built at different times and stand at different stages of
completion, but they sort by one structural feature: \emph{what moves the
rate}. Four mechanisms recur, and where a rate is driven by the process's own
state, the behaviour turns on whether the feedback exhausts a quantity or
replenishes it.

\begin{itemize}[leftmargin=1.6em,itemsep=0.35em,topsep=0.35em]
\item \textbf{Room that is spent and never returned.} Speciation under
logistic growth consumes remaining capacity $K-N(t)$, the integrated rate
converges, and the species count has a proper limit law
(\cref{var:sec:speciation}).
\item \textbf{A critical point dragged after the leader.} Resistance in the
Pimentel model acts not on a level but on the state of no preferred
direction, and so keeps restoring criticality rather than approaching an
equilibrium (\cref{var:sec:pimentel}).
\item \textbf{A potential restocked from outside.} An accumulating potential,
consumed per capita and replenished at a constant inflow, yields a cycle in
four movements (\cref{var:sec:rrrr}).
\item \textbf{A potential the population makes for itself.} Take that same
hazard and let the inflow be produced by the population's own area, with the
cost charged per \emph{type} rather than per head. The cycle becomes a
marginal state; the catastrophe divides rather than destroys; and because its
products obey the same law, the cycle branches instead of repeating
(\cref{var:sec:ca}).
\end{itemize}

The last two are a matched pair, and the chapter is arranged so that they can
be read as one movement. They also mark where the mathematics stops. The
first two sections carry closed-form results; the second two carry
simulations of models that outrun the machinery of the preceding chapters,
and say as exactly as they can what would have to be proved.

Three results carry the finished mathematics, and each is stated here so that
the derivations may be read knowing where they go.

\begin{itemize}[leftmargin=1.6em,itemsep=0.35em,topsep=0.35em]
\item For a pure birth process from one founder with per-capita rate
$\beta(t)$, the species count $S_t$ is geometric on $\{1,2,\dots\}$ with
parameter $\ee^{-\Bint(t)}$, where $\Bint(t)=\int_0^t\beta(u)\,\dd u$
(\cref{var:prop:yule-geom}). The rate law enters the distribution only
through its integral.
\item Under the logistic rate that integral converges, so the process has a
proper limit law and not merely a finite late-time mean: $S_\infty$ is
geometric with mean $(K/N_0)^{\sigma K/\varrho}$
(\cref{var:cor:logistic-limit}), depending on the parameters through exactly
those two dimensionless groups.
\item In the Pimentel competition model the state at which the chain is
equally likely to step up as down is available in closed form
(\cref{var:eq:icrit}), and it is unstable. Two trajectories started from the
same state diverge exactly when they are near it.
\end{itemize}

To these the automaton of \cref{var:sec:ca} adds an observation rather than a
theorem, and it is the one most worth taking further: a species is driven
onto the locus $X_i=\varphi N_i/c$ at which production and dissipation
cancel, and \emph{on} that locus its catastrophe hazard is proportional to
its area --- which is \ChCore{}'s own $\delta X_t$. The chapter's least
tractable model may therefore reduce to the chapter's own machinery, and
\cref{var:rem:open-ca} says what would have to be shown.

The Yule process, its geometric marginal and the method of characteristics
are \ChM{}'s; the logistic speciation model itself, with its mean, is already
a worked example there. The birth--death--catastrophe process, the roots $a$
and $b$ and the burst-size law are \ChCore{}'s and \ChDist{}'s. None of that
is re-derived here. What this chapter adds is the distribution under a moving
speciation rate, the moving critical point of the Pimentel chain, the
automaton and its balance point, and the remaining specifications, stated
precisely enough to be taken up where they stop.

\begin{table}[H]
\centering
\small
\begin{tabular}{@{}c>{\raggedright\arraybackslash}p{0.15\textwidth}>{\raggedright\arraybackslash}p{0.35\textwidth}>{\raggedright\arraybackslash}p{0.21\textwidth}@{}}
\toprule
\S & Model & What moves the rate & Status \\
\midrule
\ref{var:sec:speciation} & Logistic speciation
  & remaining capacity $K-N(t)$ & closed form \\
\ref{var:sec:pimentel} & Pimentel competition
  & evolving resistance, itself driven by the state
  & closed-form critical point; simulated \\
\ref{var:sec:pimentelplus} & Pimentel+
  & fitness driven by dominance, reduced by species count & specified \\
\ref{var:sec:rrrr} & Rise, run, ruin, rejuvenation
  & potential $V$; catastrophe by scarcity $X_t/V$ & simulated \\
\ref{var:sec:ca} & Multiplicative dissipation
  & potential $V_i$ made by area, spent per type; fragmentation at
    $\delta X_i/V_i$ & simulated \\
\midrule
\ref{var:app:plague} & Birth--plague
  & loss channel driven by the coupled component $Y_t$ & specified \\
\ref{var:app:publicgood} & Public good
  & $\lambda(t)=\lambda_0+\alpha P(t)$, $P$ made by the population
  & $K^\ast$ computed \\
\bottomrule
\end{tabular}
\caption{The models, and the mechanism by which a rate moves in each. The
last column is the shape of the chapter: \cref{var:sec:speciation} is its
mathematical end, and the rows below it halt at a closed-form fragment or at
a simulation. The two below the rule are carried as bare specifications in
\cref{var:app:plague,var:app:publicgood}, the birth--plague model in two
variants, so the count of distinct specifications is eight.}
\label{var:tab:models}
\end{table}

\subsection{Notation}
\label{var:sec:notation}

Symbols are reused across models but never within one section, and never
against a thesis-wide meaning. The table records the assignment for the body
of the chapter; the two appendix specifications gloss their own parameters
where they stand.

\begin{table}[H]
\centering
\footnotesize
\setlength{\tabcolsep}{4pt}
\begin{tabular}{@{}ll>{\raggedright\arraybackslash}p{0.56\textwidth}@{}}
\toprule
\S & Symbol & Meaning \\
\midrule
all & $\lambda,\mu,\delta$ & birth, death and catastrophe rates, as in
  \ChCore{} \\
& $a,b$ & roots of the characteristic quadratic, $b<1<a$
  (\cref{bdc:lem:roots}) \\
& $\Kb$ & burst size of a birth--death--catastrophe process \\
\midrule
\ref{var:sec:speciation}
 & $S_t$ & species count, $S_0=1$ \\
 & $N(t),K,N_0$ & total abundance, carrying capacity, initial abundance \\
 & $\varrho$ & logistic growth rate \\
 & $\sigma$ & speciation coefficient; $\beta(t)=\sigma(K-N(t))$ \\
 & $\xi$ & the constant $K/N_0-1$ \\
 & $\beta(t)$ & speciation rate per species; $\beta$ is this and nothing else \\
 & $\Bint(t)$ & integrated rate $\int_0^t\beta(u)\,\dd u$; $\eta$ is a constant floor added to $\beta$ in two of the alternative laws \\
\midrule
\ref{var:sec:pimentel},
\ref{var:sec:pimentelplus}
 & $A_t,B_t$ & counts of the two competing populations, $A_t+B_t=N$ \\
 & $N$ & capacity of the enclosure, an integer \\
 & $R_a,R_b$ & resistance of each population \\
 & $p_a,p_b$ & blocking probabilities, at most one non-zero \\
 & $\varepsilon,\alpha$ & evolution rate (the step size of the resistance update) and the exponent sharpening that drive near a boundary \\
 & $i_c$ & the critical state \cref{var:eq:icrit} \\
 & $f_a,f_b$ & fitness, the Pimentel+ recasting of $R_a,R_b$ \\
 & $S^{(a)}_t,S^{(b)}_t$ & per-population species counts; $\zeta$ couples dominance to speciation and $\gamma$ penalises fitness per species \\
\midrule
\ref{var:sec:rrrr}
 & $X_t,V(t)$ & count of replicative agencies, and accumulated potential \\
 & $\phi,c,\omega$ & potential inflow, potential consumed per agency, revival rate out of the empty state \\
\midrule
\ref{var:sec:ca}
 & $L$ & side of the lattice; $L^2$ cells, Moore neighbourhood, periodic \\
 & $N_i,X_i,V_i$ & cells held by species $i$, its extant sub-species, and
   its accumulated potential \\
 & $\varphi,c$ & potential supplied per cell, and spent per extant
   sub-species; $V_0$ is the potential a fragment is reset to \\
 & $m$ & mutation coefficient; the per-cell rate is $mV_i$ \\
 & $\theta,\nu$ & neighbour-support exponent, and the rare-type advantage
   within a species \\
 & $\kappa_i$ & mean sub-species area $N_i/X_i$; $d$ is the resistance of
   empty ground \\
\bottomrule
\end{tabular}
\caption{Notation, organised by model. The rates $\lambda$, $\mu$ and
$\delta$, the roots $a$ and $b$ and the burst size $\Kb$ carry their
thesis-wide meanings everywhere in this chapter and are never redefined.
$V$ and $c$ are shared between \cref{var:sec:rrrr,var:sec:ca} with the
\emph{same} meanings, a potential and the rate at which it is spent, which is
the point of the pairing: what differs between the two models is that
\cref{var:sec:ca} charges the inflow per cell and the outflow per type. Where
one letter genuinely serves two models --- $\alpha$ as an exponent in
\cref{var:sec:pimentel,var:sec:pimentelplus}, $N$ as an abundance in
\cref{var:sec:speciation} and an integer capacity in
\cref{var:sec:pimentel}, $X$ as a count of individuals in
\cref{var:sec:rrrr} and of types in \cref{var:sec:ca} --- the two uses are
separated by this table and never occur in the same section.}
\label{var:tab:notation}
\end{table}

\Cref{var:sec:speciation} takes the moving speciation rate as far as it will
go, which is a proper limit law. \Cref{var:sec:pimentel} obtains the moving
critical point of the Pimentel chain and, in
\cref{var:sec:pimentelplus}, specifies the coupling that would let excess
variation weaken a winner. \Cref{var:sec:rrrr,var:sec:ca} are the matched
pair: the first drives catastrophe by a potential restocked from outside, the
second by one the population makes for itself, and the second is where the
chapter's open questions concentrate. \Cref{var:app:plague,var:app:publicgood}
carry two further specifications in bare form, and \cref{var:app:pgf} gives
the generating-function proof of \cref{var:prop:yule-geom}.

\FloatBarrier
